\chapter{机器人运动学}
\label{sec:03-robot-kinematics}

\section{机器人末端位姿描述}

\begin{quote}
    Kinetic Position and Coordinate
\end{quote}

机械手的位置表达：
\begin{enumerate}
    \item 机械手末端姿态由某个姿态变换规定
    \item 相对基系的位置由\textbf{姿态矩阵}由左乘对应于矢量\(p\)的平移变换确定
\end{enumerate}

\subsection{Description in Cylinder Coordinates}

用柱面坐标表示末端位置（平移变换）

规定运动变换的过程(相对固定坐标系$xyz$运动)：
\begin{enumerate}
    \item 沿$x$轴平移$r$
    \item 绕$z$轴旋转$\alpha$
    \item 沿\(z\)轴平移\(z\)
\end{enumerate}

$$
\begin{tikzpicture}[
    % 设置 3D 透视坐标系 (Cavalier projection)
    x={(-0.5cm,-0.5cm)},
    y={(1cm,0cm)},
    z={(0cm,1cm)},
    scale=.6
]
    % 定义主要坐标点
    \coordinate (O) at (0,0,0);
    \coordinate (Px) at (3,0,0);
    \coordinate (Pxy) at (3,4,0);
    \coordinate (P) at (3,4,4.5);

    % 1. 绘制全局坐标轴 (x, y, z)
    \draw[->, thick, >=latex] (O) -- (6,0,0) node[right] {$x$};
    \draw[->, thick, >=latex] (O) -- (0,5.5,0) node[below] {$y$};
    \draw[->, thick, >=latex] (O) -- (0,0,6) node[left] {$z$};

    % 2. 绘制平面投影线及辅助线
    \draw (O) -- (Pxy); % 连向 xy 平面投影点的斜线 (r)
    \draw (Px) -- (Pxy); % 平行于 y 轴的辅助线
    \draw (Pxy) -- (P) node[midway, right] {$z$}; % 垂直于 xy 平面的高度线

    % 3. 绘制主向量
    % 使用 shorten >= 缩短一点箭头，以还原原图中箭头与局部坐标系原点间的微小空隙
    \draw[->, thick, >=latex, shorten >= 2pt] (O) -- (P);

    % 4. 绘制 xy 平面内的角度弧线和标签
    \begin{scope}[canvas is xy plane at z=0]
        % 角度 \alpha 的圆弧 (从 x 轴到投影线)
        \draw[->, >=latex] (1.5, 0) arc (0:53.13:1.5);
        \node at (26.5:1.1) {$\alpha$};
        % 标签 r 放置在靠近 x 轴的外侧
        \node at (10:2.1) {$r$};
    \end{scope}

    % 5. 绘制点 P 处的局部坐标系 (a, o, n)
    % 向量 n (指向右前/下方)
    \draw[->, thick, >=latex] (P) -- ++(2.0, 2.0, 0) node[right] {$\boldsymbol{n}$};
    % 向量 o (指向右后/上方)
    \draw[->, thick, >=latex] (P) -- ++(-1.5, 1.0, 0) node[right] {$\boldsymbol{o}$};
    % 向量 a (指向上方，平行于 z 轴)
    \draw[->, thick, >=latex] (P) -- ++(0, 0, 1.5) node[left] {$\boldsymbol{a}$};
\end{tikzpicture}
$$

变换算子：
$$
\operatorname{Cyl}(z, \alpha, r) = \operatorname{Trans}(0,0,z) \operatorname{Rot}(z, \alpha) \operatorname{Trans}(r,0,0)
$$

\begin{note}
    相对运动坐标$xyz$的变换算子：
    $$
    \operatorname{Cyl}(z, \alpha, r) = \operatorname{Trans}(0,0,z) \operatorname{Rot}(z^\prime, \alpha) \operatorname{Trans}(r^\prime,0,0)
    $$
\end{note}

\subsection{Description in Spherical Coordinates}

用球面坐标表示末端运动位置矢量

规定运动变换的过程(相对固定坐标系$xyz$运动)：
\begin{enumerate}
    \item 沿$z$轴平移$r$
    \item 绕$y$轴旋转$\beta$
    \item 沿\(z\)轴旋转\(\alpha\)
\end{enumerate}

$$
    \begin{tikzpicture}[
        >=Stealth,
        line cap=round,
        line join=round,
        scale=0.6
    ]

    %------------------------------------------------
    % Coordinates
    %------------------------------------------------
    \coordinate (O) at (0,0);
    \coordinate (X) at (-2.8,-2.0);
    \coordinate (Y) at (4.6,0);
    \coordinate (Z) at (0,5.2);

    \coordinate (B) at (2.7,-2.0);
    \coordinate (R) at (2.7,0);
    \coordinate (P) at (2.7,3.7);
    \coordinate (A) at (3.35,4.65);
    \coordinate (o2) at (4.4,4.5);
    \coordinate (N) at (4.05,2.45);

    %------------------------------------------------
    % Axes
    %------------------------------------------------
    \draw[->,thick] (O) -- (X) node[below left] {$x$};
    \draw[->,thick] (O) -- (Y) node[right] {$y$};
    \draw[->,thick] (O) -- (Z) node[above] {$z$};

    %------------------------------------------------
    % Auxiliary lines
    %------------------------------------------------
    \draw[thick] (O) -- (B);
    \draw[thick] (O) -- (R);
    \draw[thick] (B) -- (P);

    % Main vector a_0
    \draw[->,thick] (O) -- (A)
        node[pos=0.96,above right] {$\boldsymbol{a}_0$};
    % Normal vector n
    \draw[->,thick] (P) -- (N)
        node[right] {$\boldsymbol{n}$};
    % Normal vector o
    \draw[->,thick] (P) -- (o2)
        node[right] {$\boldsymbol{o}$};
    \draw[->,thick] (O) -- (P);

    %------------------------------------------------
    % Labels
    %------------------------------------------------
    \node[right] at ($(B)!0.56!(P)$) {$r$};

    %------------------------------------------------
    % Angle alpha: between x-axis and projection line
    %------------------------------------------------
    \draw[thick]
        ($(O)+(-144.5:0.75)$)
        arc[start angle=-144.5,end angle=-36.5,radius=0.75];

    \node at ($(O)+(-90:1.00)$) {$\alpha$};

    %------------------------------------------------
    % Angle beta: between z-axis and vector a_0
    %------------------------------------------------
    \draw[thick]
        ($(O)+(90:0.75)$)
        arc[start angle=90,end angle=54.2,radius=0.75];

    \node at ($(O)+(72:0.98)$) {$\beta$};

    \end{tikzpicture}
$$

变换算子：
$$
\operatorname{Sph}(z, \alpha, r) = \operatorname{Rot}(z, \alpha) \operatorname{Rot}(y,\beta) \operatorname{Trans}(0,0,r)
$$

\begin{note}
    相对运动坐标$xyz$的变换算子：
    $$
    \operatorname{Sph}(z, \alpha, r) = \operatorname{Rot}(z^\prime, \alpha) \operatorname{Rot}(y^\prime,\beta) \operatorname{Trans}(0,0,r^\prime)
    $$
\end{note}

\section{连杆参数和连杆坐标系}

\dfn{Link}{A manipulator may be thought of as a set of bodies connected in a chain by joints. These bodies are called links.}

$$
\begin{tikzpicture}[
    scale=.6, 
    > = stealth,
    parallel/.style={
        decoration={
            markings,
            mark=at position 0.6 with {
                \draw[solid, thick] (-2pt,-4pt) -- (2pt,4pt);
                \draw[solid, thick] (-6pt,-4pt) -- (-2pt,4pt);
            }
        },
        postaction={decorate}
    },
    triple parallel/.style={
        decoration={
            markings,
            mark=at position 0.6 with {
                \draw[solid, thick] (-2pt,-4pt) -- (2pt,4pt);
                \draw[solid, thick] (-7pt,-4pt) -- (-3pt,4pt);
                \draw[solid, thick] (3pt,-4pt) -- (7pt,4pt);
            }
        },
        postaction={decorate}
    }
]

    % ==========================================
    % 1. 参数与基础坐标点定义
    % ==========================================
    \def\angA{100} % 轴 i-1 的倾斜角度
    \def\angB{75}  % 轴 i 的倾斜角度
    \def\cylR{0.9} % 关节圆柱体半径
    \def\cylH{2.2} % 关节圆柱体高度

    % 公共法线交点
    \coordinate (p1) at (0,0);       % a_{i-1} 与 Axis i-1 的交点
    \coordinate (p2) at (7, -1.5);   % a_{i-1} 与 Axis i 的交点
    
    % 轴向量
    \coordinate (vAx1) at ({cos(\angA)}, {sin(\angA)});
    \coordinate (vAx2) at ({cos(\angB)}, {sin(\angB)});

    % Axis i 上的特征点
    \coordinate (p3) at ($(p2) + 1.8*(vAx2)$); % a_i 在 Axis i 上的起点 (距离 p2 上方 d_i 处)
    \coordinate (p4) at ($(p3) + (4, 0.3)$);   % a_i 延伸线的终点

    % 圆柱体底部中心点
    \coordinate (c1_bot) at ($(p1) + 2.5*(vAx1)$); % Cyl 1 (Axis i-1) 底部
    \coordinate (c2_bot) at (p3);                  % Cyl 2 (Axis i) 底部

    % 轴线的底端延伸点 (用于画线)
    \coordinate (ax1_bot) at ($(p1) - 3*(vAx1)$);
    \coordinate (ax2_bot) at ($(p2) - 3.5*(vAx2)$);


    % ==========================================
    % 2. 绘制轴线底端和法线 (a_{i-1}, a_i)
    % ==========================================
    % 轴 i-1 底部线
    \draw (p1) -- (ax1_bot);
    % 轴 i 底部线
    \draw (p2) -- (ax2_bot);
    \path[triple parallel] (p2) -- (ax2_bot); % 添加三平行线标记

    % 公共法线 a_{i-1}
    \draw[thick] (p1) -- (p2) node[midway, above=2pt] {$a_{i-1}$};
    \fill (p1) circle (1.5pt);
    \fill (p2) circle (1.5pt);

    % 法线 a_i
    \draw[thick] (p3) -- (p4) node[right] {$a_i$};
    \path[parallel] (p3) -- (p4); % 添加双平行线标记
    \fill (p3) circle (1.5pt);


    % ==========================================
    % 3. 绘制 3D 连杆 (Links)
    % ==========================================
    % --- Link i-1 (连接 Cyl 1 下部 和 Cyl 2 下部) ---
    \coordinate (L1_s1) at ($(c1_bot) + 0.3*(vAx1) + ({cos(\angA-90)*\cylR}, {sin(\angA-90)*\cylR})$);
    \coordinate (L1_s2) at ($(c1_bot) + 1.3*(vAx1) + ({cos(\angA-90)*\cylR}, {sin(\angA-90)*\cylR})$);
    
    \coordinate (L1_e1) at ($(c2_bot) + 0.1*(vAx2) + ({cos(\angB+90)*\cylR}, {sin(\angB+90)*\cylR})$);
    \coordinate (L1_e2) at ($(c2_bot) + 1.1*(vAx2) + ({cos(\angB+90)*\cylR}, {sin(\angB+90)*\cylR})$);

    \draw[thick, fill=white] (L1_s2) .. controls +(right:3.5) and +(left:3.5) .. (L1_e2) node[midway, above=8pt] {Link $i-1$}
        -- (L1_e1) .. controls +(left:3.5) and +(right:3.5) .. (L1_s1) -- cycle;

    % --- Link i (从 Cyl 2 上部向右延伸) ---
    \coordinate (L2_s1) at ($(c2_bot) + 1.2*(vAx2) + ({cos(\angB-90)*\cylR}, {sin(\angB-90)*\cylR})$);
    \coordinate (L2_s2) at ($(c2_bot) + 2.0*(vAx2) + ({cos(\angB-90)*\cylR}, {sin(\angB-90)*\cylR})$);
    
    \coordinate (L2_e2) at ($(L2_s2) + (2.5, 0.4)$);
    \coordinate (L2_e1) at ($(L2_s1) + (2.5, 0.4)$);

    \fill[white] (L2_s2) .. controls +(right:1.5) and +(left:1) .. (L2_e2) 
        -- (L2_e1) .. controls +(left:1) and +(right:1.5) .. (L2_s1) -- cycle;
    \draw[thick] (L2_s2) .. controls +(right:1.5) and +(left:1) .. (L2_e2);
    \draw[thick] (L2_s1) .. controls +(right:1.5) and +(left:1) .. (L2_e1);
    \node at ($(L2_e2) + (0, 0.5)$) {Link $i$};


    % ==========================================
    % 4. 绘制圆柱形关节 (Axis i-1 和 Axis i)
    % ==========================================
    % --- Cylinder 1 (Axis i-1) ---
    \draw (p1) -- (c1_bot); % 连接 a_{i-1} 到 Cyl 1 底部
    \begin{scope}[shift={(c1_bot)}, rotate={\angA-90}]
        % 外壳与遮挡
        \draw[thick, fill=white] (-\cylR, 0) -- (-\cylR, \cylH) arc (180:0:{\cylR} and 0.4) -- (\cylR, 0) arc (360:180:{\cylR} and 0.4);
        % 顶部椭圆与内孔
        \draw[thick] (0, \cylH) ellipse ({\cylR} and 0.4);
        \draw[thick] (0, \cylH) ellipse ({0.3*\cylR} and 0.15);
        % 轴线上半部
        \draw (0, \cylH) -- (0, \cylH+1.5) node[above] {Axis $i-1$};
    \end{scope}

    % --- Cylinder 2 (Axis i) ---
    \draw[thick] (p2) -- (p3); % d_i 段实线
    \begin{scope}[shift={(c2_bot)}, rotate={\angB-90}]
        % 外壳与遮挡
        \draw[thick, fill=white] (-\cylR, 0) -- (-\cylR, \cylH) arc (180:0:{\cylR} and 0.4) -- (\cylR, 0) arc (360:180:{\cylR} and 0.4);
        % 圆柱体中部的分割线
        \draw[thick] (-\cylR, 1.1) arc (180:360:{\cylR} and 0.4);
        % 顶部椭圆与内孔
        \draw[thick] (0, \cylH) ellipse ({\cylR} and 0.4);
        \draw[thick] (0, \cylH) ellipse ({0.3*\cylR} and 0.15);
        % 轴线上半部
        \draw (0, \cylH) -- (0, \cylH+1.5) node[above] {Axis $i$};
    \end{scope}


    % ==========================================
    % 5. 绘制直角标记 (Right Angle Symbols)
    % ==========================================
    % 直角 @ p1 (a_{i-1} 与 Axis i-1 下段)
    \coordinate (p1_a) at ($(p1)!0.3cm!(p2)$);
    \coordinate (p1_ax) at ($(p1)!0.3cm!(ax1_bot)$);
    \draw ($(p1_a)$) -- ($(p1_a) + (p1_ax) - (p1)$) -- ($(p1_ax)$);

    % 直角 @ p2 (a_{i-1} 与 Axis i 上段)
    \coordinate (p2_a) at ($(p2)!0.3cm!(p1)$);
    \coordinate (p2_ax) at ($(p2)!0.3cm!(p3)$);
    \draw ($(p2_a)$) -- ($(p2_a) + (p2_ax) - (p2)$) -- ($(p2_ax)$);

    % 直角 @ p3 (a_i 与 Axis i 下段)
    \coordinate (p3_a) at ($(p3)!0.3cm!(p4)$);
    \coordinate (p3_ax) at ($(p3)!0.3cm!(p2)$);
    \draw ($(p3_a)$) -- ($(p3_a) + (p3_ax) - (p3)$) -- ($(p3_ax)$);


    % ==========================================
    % 6. 标注 d_i 大括号
    % ==========================================
    \draw[decorate,decoration={brace,amplitude=5pt}]
        ($(p2) + ({cos(\angB+90)*0.2}, {sin(\angB+90)*0.2})$) -- 
        ($(p3) + ({cos(\angB+90)*0.2}, {sin(\angB+90)*0.2})$) 
        node[midway, left=8pt] {$d_i$};


    % ==========================================
    % 7. 标注角度 theta_i 及其平行辅助线 (使用原生 arc 重写)
    % ==========================================
    % 计算出的角度：a_{i-1}向量角度约为 -12度，a_i向量角度约为 4度
    \coordinate (theta_end1) at ($(p2) + (-12:2.5)$); 
    \coordinate (theta_end2) at ($(p2) + (4:2.5)$);  
    
    \draw[dashed] (p2) -- (theta_end1);
    \draw[dashed] (p2) -- (theta_end2);
    \path[parallel] (p2) -- (theta_end2); % 辅助线平行标记

    % 原生圆弧绘制：从 -12度 画到 4度，半径为 1.6
    \draw [->] (p2) ++(-12:1.6) arc (-12:4:1.6) node[midway, right=3pt] {$\theta_i$};


    % ==========================================
    % 8. 标注角度 alpha_{i-1} 及其平行辅助线 (使用原生 arc 重写)
    % ==========================================
    % 计算出的角度：Axis i-1 下半段角度为 -80度，Axis i 下半段角度为 -105度
    \coordinate (alpha_L1) at ($(p1) + (-80:2.5)$);
    \coordinate (alpha_L2) at ($(p1) + (-105:2.5)$);

    \draw[dashed] (p1) -- (alpha_L2);
    \path[triple parallel] (p1) -- (alpha_L2); % 辅助线三平行标记

    % 原生圆弧绘制：从 -105度 画到 -80度，半径为 1.6
    \draw [<->] (p1) ++(-105:1.6) arc (-105:-80:1.6) node[midway, below=3pt] {$\alpha_{i-1}$};
\end{tikzpicture}
$$

\begin{itemize}
    \item 杆件编号由固定基座开始，一般称基座为连杆0
    \item 关节1处于连杆1和基座之间，杆件距基座近的一端的关节为第$i$个关节
    \item 坐标系0是一个固定不动的坐标系（参考坐标系）
    \item $n$关节机器人需建立$n+1$个坐标系
\end{itemize}

\subsection{连杆参数}

两个结构参数：
\begin{itemize}
    \item 长度 \( a_{i-1} \)：第 \(i-1\) 轴指向第 \(i\) 轴的公法线长度（恒为正）。
    \item 扭角 \( \alpha_{i-1} \)：第 \(i-1\) 轴绕公垂线转至第 \(i\) 轴的夹角（可正可负）。
\end{itemize}

两个连接参数：
\begin{itemize}
    \item 偏置 \( d_i \)：两条公法线之间的距离（带正负号）。
    \item 关节角 \( \theta_i \)：两条公法线之间的夹角（带正负号）。
\end{itemize}

连杆 \(i-1\) 的四个参数为：\(a_{i-1}\)、\(\alpha_{i-1}\)、\(d_i\)、\(\theta_i\)。

\begin{note}
    \begin{itemize}
        \item \textbf{旋转关节 (Revolute Joint)}：关节变量 \(\theta_i\)，连杆参数 \(a_{i-1}\)、\(\alpha_{i-1}\)、\(d_i\) 不变。
        \item \textbf{移动关节 (Prismatic Joint)}：关节变量 \(d_i\)，连杆参数 \(a_{i-1}\)、\(\alpha_{i-1}\)、\(\theta_i\) 不变。
    \end{itemize}
\end{note}

\subsection{Denavit-Hartenberg方法}

连杆$i-1$的四个参数：$a_{i-1}$、$\alpha_{i-1}$、$d_i$、$\theta_i$

\begin{itemize}
    \item 旋转关节 (Revolute Joint)：关节变量 \( \theta_i \)，连杆参数 \( a_{i-1} \)、\( \alpha_i \)、\( d_i \) 不变；
    \item 移动关节 (Prismatic Joint)：关节变量 \( d_i \)，连杆参数 \( a_{i-1} \)、\( \alpha_i \)、\( \theta_i \) 不变；
    \item 六自由度机器人：18个参数表示其运动学参数（固定），6个关节变量是运动学方程中的变量。
\end{itemize}



\dfn{Denavit—Hartenberg notation}{Hence, any robot can be described kinematically by giving the values of four
quantities for each link. Two describe the link itself, and two describe the link's
connection to a neighboring link. In the usual case of a revolute joint, is called
the joint variable, and the other three quantities would be fixed link parameters.
For prismatic joints, d1 is the joint variable, and the other three quantities are fixed
link parameters. The definition of mechanisms by means of these quantities is a
convention usually called the Denavit—Hartenberg notation}

0号连杆（基座）的参数定义:
\begin{itemize}
    \item 0号连杆为基座，没有前向关节
    \item $a_0=0$、 $\alpha_0=0$
    \item 如果1号关节为旋转关节，则$d_1=0$
    \item 如果1号关节为移动关节，则$\theta_1=0$
\end{itemize}

$n$号连杆（末端）的参数定义
 \begin{itemize}
    \item 末端没有后向关节
    \item $a_n=0$、 $\alpha_n=0$
    \item 如果末端关节为旋转关节，则$d_n=0$
    \item 如果末端关节为移动关节，则$\theta_n=0$
 \end{itemize}

\ex{连杆坐标系建立的步骤}{
    \begin{enumerate}
        \item 找出各个关节轴
        \item 画出相邻关节轴线的公垂线
        \item 规定$z_i$轴与关节轴$i$重合
        \item 规定$x_i$轴与公垂线$a_i$重合
        \item 定义$y_i = z_i \times x_i$
        \item 按规则确定$\{0\}$、$\{n\}$
    \end{enumerate}
}

\subsection{Craig法则 (Modified DH)}

\begin{figure}[H]
    \centering
    \includegraphics[width=.8\textwidth]{Craig法则.png}
\end{figure}

中间连杆的坐标系$\{i-1\}$：
\begin{itemize}
    \item $z$轴：与关节轴线 \(i\) 共线，指向任意。
    \item $x$轴：与连杆公法线 \(a_i\) 重合，指向由关节i到关节 \(i+1\)；当 \(a_i = 0\) 时，取 \(x_i = \pm z_{i+1} \times z_i\)。
    \item $y$轴：取 \(y_i = z_i \times x_i\)（按右手法则）。
    \item 原点$O$：取在 \(x_i\) 和 \(z_i\) 的交点上。
        \begin{itemize}
            \item 当 \(z_{i-1}\) 与 \(z_i\) 相交，取两轴交点使 \(a_{i-1} = 0\)。
            \item 当 \(z_{i-1}\) 与 \(z_i\) 平行，取使 \(d_{i-1} = 0\)。
        \end{itemize}
\end{itemize}

\begin{minipage}{0.45\textwidth}
    \centering
    \includegraphics[width=\linewidth]{连杆坐标系定义——基座.png} % 图片宽度占满左栏
\end{minipage}
\hfill % 水平填充，拉开左右距离
\begin{minipage}{0.45\textwidth}
    连杆坐标系定义——基座
    \begin{itemize}
        \item 基座是连杆0，没有前向关节
        \item 规定坐标系$\{0\}$原点与坐标系$\{1\}$原点重合
        \item $\hat{Z}_0$轴与$\hat{Z}_1$轴重合
        \item 当关节变量$q_1=0$时，$\hat{X}_0$轴与$\hat{X}_1$轴重合
        \item $\hat{Y}_0$轴由右手定则确定
        \item 坐标系$\{0\}$是固定坐标系
    \end{itemize}
\end{minipage}

\vspace{1cm}

\begin{minipage}{0.45\textwidth}
    \centering
    \includegraphics[width=\linewidth]{连杆坐标系定义——末端连杆.png} % 图片宽度占满左栏
\end{minipage}
\hfill % 水平填充，拉开左右距离
\begin{minipage}{0.45\textwidth}
    连杆坐标系定义——末端连杆
    \begin{itemize}
        \item 末端连杆$n$没有后向关节
        \item 规定坐标系$\{n\}$的原点在$a_{n-1}$与关节$n$轴线$S_n$的交点处
        \item $\hat{Z}_0$轴与关节$n$轴线$S_n$轴重合
        \item 当关节变量$q_n=0$时，$\hat{X}_n$轴与$\hat{X}_{n-1}$轴共线且同向
        \item $\hat{Y}_0$轴由右手定则确定
    \end{itemize}
\end{minipage}

\begin{note}
    连杆坐标系定义小结——步骤
    \begin{enumerate}
        \item 连杆和关节编号规定：从基座到末端依次编号，基座为连杆0，连接基座与连杆1的关节为关节1；
        \item 画出各关节轴线，标记轴线$i$和$i+1$的交点，或两者的公垂线$a_i$及其与轴$i$的交点，该交点为坐标系$\{i\}$的原点；
        \item 确定$\hat{Z}_i$轴沿轴$i$的方向（有两种选择）；
        \item 规定$\hat{Z}_i$轴沿公垂线$a_i$从轴i指向轴$i+1$。若轴$i$与轴$i+1$相交，则$\hat{Z}_i$轴垂直于轴$i$与轴$i+1$所在平面，此时$\hat{Z}_i$轴的方向有两种选择；
        \item 按照右手定则确定$\hat{Y}_i$轴
        \item 坐标系总数量为$N+1$个，$N$为关节数；
        \item 对于基座，规定其坐标系编号为$\{0\}$，与坐标系$\{1\}$重合；对于末端坐标系$\{N\}$，其原点在$a_{n-1}$与轴N的交点上，当$q_n=0$时，$\hat{X}_n$轴与$\hat{X}_{n-1}$轴共线且方向相同。
    \end{enumerate}
    \begin{figure}[H]
        \centering
        \includegraphics[width=.5\textwidth]{连杆坐标系定义小结——步骤.png}
    \end{figure}
\end{note}

\subsection{坐标系变换}

连杆坐标系$\{i-1\}$到坐标系$\{i\}$的齐次变换：
$$
{}^{i-1}_i T
$$
步骤：
\begin{enumerate}
    \item $\{i-1\}$绕$\hat{X}_{i-1}$旋转$\alpha_{i-1}$，得到$\{R\}$
    \item $\{R\}$$\hat{X}_{i-1}$移动$a_{i-1}$，得到$\{Q\}$
    \item $\{Q\}$绕$\hat{Z}_i$旋转$\theta_i$，得到$\{P\}$
    \item $\{P\}$沿$\hat{Z}_i$平移$d_i$，与$\{i\}$重合
\end{enumerate}

\[
{}^{i-1}_i T = {}^{i-1}_R T {}^R_Q T {}^Q_P T {}^P_i T = \operatorname{Rot}(x, \alpha_{i-1}) \operatorname{Trans}(x, a_{i-1}) \operatorname{Rot}(z, \theta_i) \operatorname{Trans}(z, d_i)
\]

连杆坐标系变换过程
\begin{enumerate}
    \item 绕\(x_{i-1}\)轴转\(\alpha_{i-1}\)角;
    \item 沿\(x_{i-1}\)轴移动\(a_{i-1}\);
    \item 绕\(z_{i}\)轴转\(\theta_{i}\)角;
    \item 沿\(z_{i}\)轴移动\(d_{i}\).
\end{enumerate}

\subsection{广义变换矩阵}

相邻坐标系间的关系
\[
\begin{aligned}
    {}^{i-1}_i T &=\operatorname{Rot}(x, \alpha_{i-1}) \operatorname{Trans}(x, a_{i-1}) \operatorname{Rot}(z, \theta_i) \operatorname{Trans}(z, d_i) \\
    &=\begin{bmatrix} \cos\theta_i & -\sin\theta_i & 0 & a_{i-1} \\ \sin\theta_i\cos\alpha_{i-1} & \cos\theta_i\cos\alpha_{i-1} & -\sin\alpha_{i-1} & -d_i\sin\alpha_{i-1} \\ \sin\theta_i\sin\alpha_{i-1} & \cos\theta_i\sin\alpha_{i-1} & \cos\alpha_{i-1} & d_i\cos\alpha_{i-1} \\ 0 & 0 & 0 & 1 \end{bmatrix}
\end{aligned}
\]

\begin{note}
    平移矩阵
    $$
    \operatorname{Trans}({}^A\boldsymbol{p}) = \begin{bmatrix} 
    1 & 0 & 0 & {}^A p_x \\ 
    0 & 1 & 0 & {}^A p_y \\ 
    0 & 0 & 1 & {}^A p_z \\ 
    0 & 0 & 0 & 1 
    \end{bmatrix}
    $$
    旋转矩阵
    $$
    \boldsymbol{R}(\boldsymbol{k},\theta) = \begin{bmatrix} 
    \cos\theta & -\sin\theta & 0 & 0 \\ 
    \sin\theta & \cos\theta & 0 & 0 \\ 
    0 & 0 & 1 & 0 \\ 
    0 & 0 & 0 & 1 
    \end{bmatrix}
    $$
\end{note}

末端$\{n\}$相对基坐标系$\{0\}$的关系
$$
{}^{0}_n T\left(\theta_1, \theta_2, \theta_3, \cdots, \theta_n\right) = {}^{0}_1 T {}^{1}_2 T \cdots {}^{n-1}_n T
$$

运动学方程：表示 $\underset{(\boldsymbol{n}, \boldsymbol{o}, \boldsymbol{a}, \boldsymbol{p})}{\text{\textbf{末端位姿}}}$ 相对于基座的位置关系，是 $\underset{\theta_1,\dots,\theta_n}{\text{\textbf{关节变量}}}$ $(a, \alpha, d, \theta)$ 的函数。
